% !TEX root = ../notes.tex

\documentclass[../notes.tex]{subfiles}

\begin{document}

\section{March 9}
We will assume that we know something about reductive groups already.

\subsection{Historical Remarks}
Let's discuss the history of the following theorem.
\begin{theorem}
	For a reductive group $G$, there are only finitely many unipotent conjugacy classes.
\end{theorem}
Reductive groups go back to Killing, who apparently did not know that symplectic groups exist. Eventually, he managed to construct $G_2$. Conjugacy classes (in the form of the Jordan normal theorem) go back to Weierstrass, though often attributed to Jordan. In particular, by the Jordan normal theorem, conjugacy classes are parameterized by a pair of a unipotent conjugacy class and a semisimple conjugacy class. Over $\CC$, semisimple conjugacy classes is not so hard (they all diagonalize). For unipotent conjugacy classes, the classification is due to Weierstrass for $\mf{gl}_n$ (via some partitions) and due to Williamson for orthogonal groups.

Let's discuss the finiteness of unipotent classes. The first important step was by Morosov, who showed that unipotent classes in a reductive group $G$ are parameterized by algebraic maps from $\op{SL}_2$. The Myatev showed that there are only finitely many such maps up to conjugacy, which together shows that there are only finitely many unipotent conjugacy classes. This result is frequently attributed to Dynkin, who shortly fully classifies all such maps up to conjugacy. Later, it was shown that the parameterization given by Morosov was in fact a bijection, therefore enumerating conjugacy classes.

Chevalley showed that the theory of reductive groups works well over any algebraically closed field, but the discussion of the previous paragraph does not generalize to positive characteristic (even for good primes): Morosov's theorem is no longer true.
\begin{example}
	In $E_8$, there are $70$ conjugacy classes in characteristic zero, $74$ conjugacy classes in characteristic two, and $71$ conjugacy classes in characteristic three.
\end{example}
\begin{remark}
	In general, there is a natural embedding from conjugacy classes in characteristic zero to conjugacy classes in positive characteristic. However, there may be more conjugacy classes in the latter!
\end{remark}
Wall would classify conjugacy classes for classical groups in characteristic $2$. For good primes, the finiteness was proved for exceptional groups (though without classification). Students of Iwahori would classify these conjugacy classes shortly. Lusztig (based on work with Deligne) showed finiteness in generality shortly thereafter, which was new for the exceptional types $E$, and classification followed in types $E$ by Mizuno.

\subsection{Characteristic Zero to Positive Characteristic}
Let's begin doing math.
\begin{notation}
	Fix a connected reductive group $G$ over an algebraically closed field $k$. Then we define $\mc D_G$ to be the set of algebraic group homomorphisms $\omega\colon\mathbb G_m\to G$ such that $\omega$ lifts to an algebraic homomorphism $\widetilde\omega\colon\mathrm{SL}_2\to G$ for which
	\[\widetilde\omega\left(\op{diag}(t,1/t)\right)=\omega(t).\]
	Given some $\omega\in\mc D_G$, we define a descending filtration $\{G_n^\omega\}$ of $G$ so that $\op{Lie}G_n^\omega=\mf g_n\oplus\mf g_{n+1}\oplus\mf g_{n+2}\oplus\cdots$, where $\mf g_n\coloneqq\{X\in\mf g:\omega(t)=t^nX\}$.
\end{notation}
\begin{remark}
	It turns out that $G_0^\omega$ is parabolic, and $G_1^\omega$ is its unipotent radical. (These claims are checked by passing to Lie algebras. For example, for the former claim, one needs to show that $\mf g_{\ge n}^\omega$ contains a Borel subgroup.) Furthermore, $G_2^\omega$ is contained in the unipotent elements of $G$.
\end{remark}
\begin{remark}
	We work over $k=\CC$. The natural map from $\op{Hom}(\mathrm{SL}_2,G)$ to unipotent classes (which descends to a bijection up to conjugacy) factors through $\mc D_G$ simply by restricting the map $\mathrm{SL}_2\to G$ to the diagonal. One recovers a map from $\mc D_G$ to the unipotent classes by looking at the elements in $G_2^\omega$.
\end{remark}
\begin{notation}
	Fix a connected reductive group $G$ over an algebraically closed field $k$. For a $G$-orbit $\overline\omega$ of some $\omega\in\mc D_G$, we define
	\[\widetilde H^{\overline\omega}\coloneqq\bigcup_{\omega\in\overline\omega}G_2^\omega,\]
	and
	\[H^{\overline\omega}\coloneqq\widetilde H^{\overline\omega}\setminus\bigcup_{\substack{\overline\omega'\in\mc D_G/G\\\overline\omega'\ne\overline\omega}}\widetilde H^{\overline\omega'}.\]
	These $H^{\overline\omega}$s turn out to partition the unipotent variety.
\end{notation}
\begin{remark}
	The set $\mc D_G$ is independent of the field (it only depends on root datum). The $H^{\overline\omega}$ contains a unique unipotent class, so the moral is that there is an embedding from unipotent classes of $G$ over $\CC$ to unipotent classes of $G$ over $k$.
\end{remark}

\subsection{The Springer Resolution}
We will want to know a little bit about the Springer resolution.
\begin{notation}
	Fix a connected reductive group $G$ over an algebraically closed field $k$. Then $\mc B$ is the flag variety of Borel subgroups, and $W$ is the Weyl group.
\end{notation}
\begin{remark}
	The Bruhat decomposition allows us to parameterize $W$ as $G$-orbits in $\mc B\times\mc B$. 
\end{remark}
\begin{definition}[Springer resolution]
	Fix a connected reductive group $G$ over an algebraically closed field $k$. Then we have a \textit{Springer--Grothendieck resolution}
	\[\pi\colon\widetilde{\mc B}\to\mc B,\]
	where $\widetilde{\mc B}\coloneqq\{(g,B)\in G\times\mc B:g\in\mc b\}$. We let $\widetilde{\mc B}_{\mathrm{rs}}$ denote the pre-image of the regular semisimple locus of $G$.
\end{definition}
\begin{remark}
	The Springer resolution is generically $W$-to-$1$. Indeed, $\widetilde{\mc B}_{\mathrm{rs}}\to\mc B$ is an \'etale covering with fibers given by $W$. Thus, one has an embedding $W\to\op{Aut}\widetilde{\mc B}_{\mathrm{rs}}$, which turns out to be an embedding of groups.
\end{remark}
This map $\pi$ has some special properties.
\begin{theorem}[Lusztig]
	The map $\pi\colon\widetilde{\mc B}\to\mc B$ is small, meaning that the dimensions of the fibers are small in some technical sense.
\end{theorem}
\begin{remark}
	Formal nonsense takes the smallness of $\pi$ and produces an action of $W$ on $\mathrm H^*(\mc B_g;\overline\QQ_\ell)$ using some intersection cohomology. (Here, $\mc B_g$ is the fiber of $\mc B$ over $g$.) Such an action was expected over $\CC$, but this definition now works in any characteristic.
\end{remark}
We are now ready to produce our Springer representation.
\begin{notation}
	Fix a connected reductive group $G$ over an algebraically closed field. For $g\in G$, we let $V_g$ be the co-image of the map
	\[\mathrm H^*(\mc B;\overline\QQ_\ell)\to\mathrm H^*(\mc B_g;\overline\QQ_\ell)\]
	in degree $2\dim\mc B_g$.
\end{notation}
\begin{remark}
	The action of $W$ on $\mc B$ produces an action on $\mathrm H^*(\mc B;\overline\QQ_\ell)$. Thus, $V_g$ is a representation of $W$, and it turns out to be irreducible. In terms of homology, it is equivalent to ask for the image of this map (in top-degree of the right-hand homology group) to be irreducible.
\end{remark}

\subsection{Our Strata}
Here is the main definition of our course.
\begin{notation}
	Fix a connected reductive group $G$ over an algebraically closed field. Then $g,g'\in G$ are \textit{equivalent}, written $g\sim g'$, if and only if $V_g\cong V_{g'}$ as representations of $W$. This stratifies $G$.
\end{notation}
We will find these strata very interesting. Later, we will show the following.
\begin{itemize}
	\item We will show that each stratum is a union of conjugacy classes of fixed dimension.
	\item Each stratum contains at most one unipotent conjugacy class.
	\item The strata are in bijection with the irreducible representations of $W$.
	\item Any stratum contains a unipotent element for some choice of characteristic.
	\item The strata turn out to be the same (roughly speaking) for a group and its Langlands dual. Notably, this is not true for unipotent classes!
\end{itemize}
\begin{example}
	There are three unipotent classes for $G\coloneqq\op{GL}_3$ over $\CC$, and it turns out that there are three strata. Let $G_d\subseteq G$ be the subset of those $g\in G$ for which $\dim Z_G(g)=d$. It turns out that $d\in\{0,4,6\}$, and these sets turn out to be the strata. (On the homework, we will check that $\dim Z_G(g)\in\{0,4,6\}$, and we will check that the $G_d$s are smooth irreducible varieties.)
\end{example}
\begin{example}
	We work with $G\coloneqq\op{Sp}_4$ over $\CC$, and again let $G_d\subseteq G$ be the subset of those $g\in G$ for which $\dim Z_G(g)=d$. The $G_d$s turn out to not be the strata: indeed, $G_4$ splits into $G_4'\sqcup G_4''$, where the former is the conjugacy class of $\op{diag}(1,1,-1,-1)$, and the latter comes from a unipotent conjugacy class. But up to this splitting, we have indeed recovered our strata. Notably, $G_4'$ contains no unipotent class in characteristic zero, but it does have a conjugacy class in characteristic two. We are asked to show that these strata have smooth irreducible components on the homework, but it may be too difficult.
\end{example}
\begin{remark}
	In general, the strata have smooth irreducible components for classical groups, but Professor Lusztig does not know if this is true in general.
\end{remark}

\end{document}